\section{Attack Asymmetry、Shared Signals 与 Zero Trust}

攻击者只需找到一个 credential、session、help-desk 或 integration gap，防守方却要协调 endpoint、identity、application、data 和 human workflow。课堂把这种不对称与信息共享障碍联系起来；系统设计需要 defense in depth，也需要在隐私和法律边界内交换 risk signal，减少每家组织独立发现同一攻击模式的时间。

\subsection{一处弱点对多层防线}

防守团队不能把所有风险压给 IdP。Endpoint compromise、phishing、malware、OAuth consent、stale account、over-privileged service account 和 application bug 都可能绕过单一检查。Identity layer 的优势是看到跨应用登录、factor、session 和 group 状态，但它仍需接收 EDR、network、application 与 customer-reported indicators。

\lecturefigure{10-shared-defense.png}{攻击者需要一个缺口，防守方要协调多层控制与跨组织信号。}{本地字幕 21:20--24:10；概念重绘。}

\begin{knowledgebox}{读图：Shared signal 也要最小化}
Indicator 可以是 suspicious IP、session risk、credential compromise 或 device posture；共享前应定义 schema、confidence、purpose、retention 和 recipient。传递过少会失去防御价值，传递过多则扩大隐私和商业敏感数据风险。Signal 应触发本地 policy evaluation，而不是让外部来源直接拥有最终 deny 权限。
\end{knowledgebox}

\subsection{Zero Trust：显式、上下文与可撤销的信任}

\term{Zero Trust} 不是“永远不信任任何人”，而是不因网络位置或一次登录永久授予信任。NIST SP 800-207 强调每次访问按 subject、asset、environment 与 policy 做显式决策，采用 least privilege，并假设网络可能已被攻破。现代 identity policy 还会在 session 中持续接收风险信号并 step-up、降权或 logout。

\lecturefigure{11-zero-trust-policy.png}{Zero Trust 把身份、设备、风险与资源上下文转成持续可撤销的 policy。}{NIST SP 800-207；Okta policy/Identity Threat Protection；概念重绘。}

\begin{knowledgebox}{读图：Policy evaluation 不止发生在登录时}
Identify 建立 principal 与 authenticator；evaluate 汇集 device、network 和 risk；enforce 输出 least privilege、step-up 或 deny；re-evaluate 在 session 风险变化后 revoke 或 universal logout。持续评估并不意味着无限监控，而是为高价值 action 选择必要信号，并让 session 拥有明确 expiry 与 revocation path。
\end{knowledgebox}

\begin{warningbox}{Zero Trust 不能修复混乱的 entitlement}
如果组织不知道谁拥有 service account、哪些 group 对应何种权限、离职何时传播或应用是否支持 revoke，那么更频繁认证只能反复验证同一个过度授权主体。Zero Trust 前提是 inventory、ownership、least privilege 和 lifecycle hygiene。
\end{warningbox}

\subsection{本章小结}

Attack asymmetry 要求多层防御和共享信号，Zero Trust 则把信号转成显式、最小权限、可撤销的访问决策。两者都依赖清洁 identity state 与审计边界。

